% SHARD-ID: CA-90-LOCAL-UNITAL-LIFT
% SHARD-TITLE: Local Unital Lifts from Relative Half-Braidings
% SHARD-SUMMARY: For unitary fusion C and pointed O = 1 + X, a pointed unitary relative half-braiding of O on <O> (a braiding suffices) gives the dyadic lift alpha_L(a) = B_L(a (x) 1)B_L^dagger: a unital injective *-map extending vacuum insertion, support [2j-1,2k-1] (r = 0), exact occupation degree, natural, WT-1 A1; CA-76's strict drift bound applies, A3 does not follow.
% SHARD-SUMMARY: Finite half-braiding equations with the integer condition [O][x] = [x][O] in K(<O>), which (review Prop. R-1) obstructs every one-gap lift of the same support pattern; the contiguous shuffle needs depth L-1 (optimal). Closes CA-75 Link (b1) and supersedes Proposal 76.5 at this layer.
% SHARD-KEYWORDS: local unital lift, relative half-braiding, Drinfeld centre, pointed site object, vacuum insertion, dyadic refinement, WT-1, A1, Wilson drift, occupation grading, Grothendieck ring, one-gap lift, Bratteli inclusion matrix, circuit depth, Link (b1), Proposal 76.5, R0^u

\section{Local Unital Lifts from Relative Half-Braidings}
\label{sec:local-unital-lift}

\subsection{Status, Sources, and Dependencies}

\textbf{Status: mixed; every statement carries its own label.}  Derived: Lemma 90.0, Theorem 90.1, Negative 90.N, Corollary
90.W, Proposition 90.V, Proposition 90.2 (part (b) is the reviewer's Proposition R-1), Proposition 90.3, Corollary 90.D (an
implication from a \emph{supplied} circuit).  Proposal: Target 90.T.  Open: arbitrary local dyadic lifts at \(\mathsf L0\),
WT-1 A3 for any non-trivial state family, the periodic seam, parity covariance.  Review grades (carried, not upgraded):
braided headline \textbf{PROOF-COMPLETE}; non-braided converse \textbf{WEAKER-RESULT-PROVED}, sharpened by R-1 (E1
REVIEW:12--21).  \textbf{Abbreviations:} ``E1 NOTE/REVIEW/INTEGRATION'' are \path{NOTE.md}, \path{REVIEW.md},
\path{INTEGRATION.md} in \path{explorations/2026-10-22-astra-r2/E1-local-unital-lift/}, whose \path{checks/} hold Python/NumPy
Float64 calibrations (commands: \path{checks/README.md}:8--13, REVIEW:180--183) that calibrate and do not prove; a port to
\path{scripts/} with a run bundle is open.  No Julia or \LaTeX{} run was made.  Proofs are structured, not proof-assistant.
This shard \textbf{closes CA-75 Link (b1)} and \textbf{supersedes CA-76 Proposal 76.5} for braided \(\Cc\) and for
\(\mathsf L0\) with a declared pointed relative half-braiding, and answers D3 open problem 1 at that layer.
Dependencies: CA-68, CA-71, CA-75, CA-76, CA-88; WT-1; CONVENTIONS (r), (t), (x), (am), (an), (aq).

\textbf{Layer tags [(aq)].}  Setting is \(\mathsf{any}\): no target theory enters; the microscopic declaration is an
\emph{open} chain with declared endpoints (REVIEW:114--119).  ``\(+\)H'' means Definition 90.H's datum.
\begin{center}\footnotesize
\begin{tabular}{p{4.4cm}p{1.0cm}p{4.4cm}p{4.2cm}}
\toprule
Statement & \(\mathsf S\) & \(\mathsf L\) & \(\mathsf R\) \\
\midrule
Lemma 90.0, Theorem 90.1, Neg.\ 90.N & any & \(\mathsf L0\)+H (in particular \(\mathsf L1\)) & \(\mathsf R0^{u}\) \\
Corollary 90.W & any & \(\mathsf L0\)+H, declared states & \(\to\mathsf R1\) (conditional on summability) \\
Proposition 90.V & any & \(\mathsf L0\)+H, constructed states & \(\mathsf R1\) (vacuum product only); negative \\
Proposition 90.2 (a); (b) & any & \(\mathsf L0\) & \(\to\mathsf R0^{u}\); negative for one-gap lifts \\
Proposition 90.3, Corollary 90.D & any & \(\mathsf L0\)+H (90.D: \(+\) a supplied circuit) & \(\mathsf R0^{u}\) implementation \\
Target 90.T & any & \(\mathsf L0\) & \(\to\mathsf R0^{u}\) [proposal] \\
\bottomrule
\end{tabular}
\end{center}

% Source: explorations/2026-10-22-astra-r2/E1-local-unital-lift/NOTE.md:19-26 (H); :30-68 E1.1; :72-117 E1.2; :119 absolute
%   number; :123-147 E1.3-E1.4; :151-160 E1.5; :162-191 E1.6, boundary; :195-224 E1.7, D3 dressing; :234-284 ledger, tags, examples.
% Source: .../E1-local-unital-lift/REVIEW.md:10-21 verdict; :23-43 citations; :45-120 step verdicts; :122-146 refutation
%   attempts; :148-168 Proposition R-1; :170-176 corrections.  INTEGRATION.md:6-31 statements; :49-91 amendments, links, next lemma.
% Source: explorations/2026-10-22-astra-r2/briefs/E1-local-unital-lift.md:50-81 -- Conjecture E1 and the converse question.
% Source: report/sections/68_variable_n_refinement_maps.tex:51-92 (V_L, corner, Boundary); 71_...:102-113 (V_psi V_phi = V_psi.phi);
%   76_...:75-81 (76.B), :90-108 (76.1), :133-147 (76.4, 76.5), :236-247 (76.12); 75_integration_map.tex:149-155 (Link (b1)).
% Source: references/text/PenneysUnitaryFusionCategories.md:219-237, :319, :325-348, :395-398 -- splittings, unitary tensor
%   structure, monoidal functors/tensorators, associator suppression.  FuchsRunkelSchweigert2002/source/hep.tex:806-819, :921-924,
%   :1023-1125 (references/categorical-reconstruction/) -- coherence, braiding naturality and tensoriality.
% Source: literature/md/1704.04221/1704.04221.md:298-307 (half-braiding); references/text/EtingofNikshychOstrikFusionCategories.md:1221
%   (K(Z(C)) -> centre of K(C)); cft_machine/research/wilson_triangle.md:18-129 (WT-1, WT-2); literature/md/2010.11121/2010.11121.md:164-176.
% Run: .../E1-local-unital-lift/checks/output.txt:1-6; red_phase_free.txt:12; review_pointed_lift_output.txt:1-13;
%   review_one_gap_obstruction_output.txt:1-4.

\subsection{Setting and the Transport Datum}

\(\Cc\) is unitary fusion and \((O,\eta)\) is pointed as in (am): \(\eta^\dagger\eta=1\), \(p_0=\eta\eta^\dagger\),
\(p_1=1-p_0\), \(O\cong\one\oplus X\); \(X\) may contain unit summands and need not be simple, self-dual or
multiplicity-free.  \(W_M=O^{\otimes M}\) (left-associated, (r)), \(\Alg_M=\operatorname{End}_\Cc(W_M)\) acting on
\(\bigoplus_c\Hom(c,W_M)\), and \(\mathcal D=\langle O\rangle\) is the \emph{full} fusion subcategory generated by \(O\), so
every \(W_M\) and every endomorphism of it lies in \(\mathcal D\).  \textbf{Support} means membership:
\(a\in\Alg_{[j,k]}:=1_{W_{j-1}}\otimes\operatorname{End}(W_{k-j+1})\otimes1_{W_{M-k}}\) (single-site algebras need not
generate \(\Alg_M\)).  \textbf{Degree}: with \(N_M=\sum_jp_1^{(j)}\), \(a\) has degree \(q\) iff \([N_M,a]=qa\); this is the
(am) grading since the \(Q_{M,n}\) are the spectral projections of \(N_M\) (REVIEW:49--53).  \(V_L\) is the CA-68 dyadic
insertion (placement \(\varphi(j)=2j-1\), (t)), \(\theta_L=\mathrm{Ad}\,V_L\) its corner map.

\textbf{Definition 90.H [definition; \(\mathsf L0\) plus a declared datum; NOTE:19--26].}  A \emph{pointed unitary
relative half-braiding} of \((O,\eta)\) on \(\mathcal D\) is a family of unitaries \(t_Y:Y\otimes O\to O\otimes Y\),
\(Y\in\mathcal D\), with (associators and unitors suppressed; Penneys:319, :395--398; FRS:806--819)
\[
 t_\one=1_O,\quad t_Z(f\otimes1_O)=(1_O\otimes f)t_Y\ (f:Y\to Z),\quad t_{Y\otimes Z}=(t_Y\otimes1_Z)(1_Y\otimes t_Z),\quad
 t_Y(1_Y\otimes\eta)=\eta\otimes1_Y. \tag{H}
\]
\(z_Y=t_Y^{-1}\) is a half-braiding of \(O\) relative to \(\mathcal D\) (1704.04221:298--307); unitarity and pointedness are
imposed, not inferred from that source.

\textbf{Remark 90.R [derived; NOTE:16, :160, :283--284; REVIEW:49--53, :79--83].}  (a) A unitary braiding \(c\) on
\(\Cc\), or only on \(\mathcal D\), gives \(t_Y=c_{Y,O}\) satisfying (H) (FRS:921--924, :1023--1125); the \(\eta\) identity
uses \(c_{Y,\one}=\mathrm{id}\), standard but not displayed at a cited local line (\emph{Law 1 gap}, REVIEW:31).  The
opposite crossing gives a different valid tower; no mixture is made.  (b) A lift of \(O\) to \(Z(\Cc)\) with a central
pointed inclusion restricts to (H) but is strictly stronger (\(\mathrm{Vec}_G\) row of the table).  (c) Centrality of the
chosen \(\eta\) is needed: a half-braiding on \(O\) with non-central point breaks intertwining (\(\mathrm{Vec}_{\mathbb Z/2}\)
row).  (d) Conversely a central \(\eta\) splits \(O=\one\oplus X\) in \(Z(\mathcal D)\), as \(p_0,p_1\) are then central.

\subsection{The One-Gap Lift}

\textbf{Lemma 90.0 [derived; E1.1, NOTE:30--68; REVIEW:55--61, every step valid].}  \textsc{Assume} (H), \(0\le s\le M\),
\(b_{M,s}=1_{W_s}\otimes t_{W_{M-s}}:W_M\otimes O\to W_{M+1}\) (new site after old site \(s\)),
\(\beta_{M,s}(a)=b(a\otimes1_O)b^\dagger\), and \(v_{M,s}\) the \(\eta\)-insertion in that gap.  \textsc{Prove}: \(\beta\) is a
unital injective \(*\)-homomorphism, \(\beta(a)v=va\), and for \(a\in\Alg_{[j,k]}\)
\[
 \beta(a)\in\Alg_{[j,k]}\ (s\ge k),\qquad \Alg_{[j+1,k+1]}\ (s<j),\qquad \Alg_{[j,k+1]}\ (j\le s<k). \tag{90.1}
\]
\(\langle1\rangle1\). \(b\) is unitary; tensor interchange and dagger (Penneys:319) give a unital \(*\)-map; compressing by
\(1\otimes\eta\) returns \(a\), so it is injective.
\(\langle1\rangle2\). The \(\eta\) identity on the suffix gives \(b(1\otimes\eta)=v\), so \(\beta(a)v=b(1\otimes\eta)a=va\).
\(\langle1\rangle3\). \(s\ge k\): \(b\) acts only after site \(k\) and on the ancilla; interchange.
\(\langle1\rangle4\). \(s<j\): \(a=1_{W_s}\otimes a_{\rm tail}\) with \(a_{\rm tail}\) an endomorphism of the whole suffix;
naturality gives \(t(a_{\rm tail}\otimes1)t^\dagger=1_O\otimes a_{\rm tail}\).  No generation by one-site operators is used;
this is where (H) on the \emph{full} \(\mathcal D\) is needed.
\(\langle1\rangle5\). \(j\le s<k\): with \(A=W_{j-1}\), \(J=W_{s-j+1}\), \(R=W_{k-s}\), \(C'=W_{M-k}\), the tensor identity gives
\(b=(1_{AJ}\otimes t_R\otimes1_{C'})(1_{AJR}\otimes t_{C'})\); the second factor commutes past \(a\), leaving
\(1_A\otimes[(1_J\otimes t_R)(a_{JR}\otimes1_O)(1_J\otimes t_R^\dagger)]\otimes1_{C'}\in\Alg_{[j,k+1]}\).  \(\square\)
Only (H) is used: no self-duality, multiplicity-freeness, ribbon structure or TL quotient (REVIEW:61).

\subsection{The Dyadic Lift}

\textbf{Theorem 90.1 [derived; E1.2 (braided) and E1.5 (\(\mathsf L0\)+H), NOTE:72--117, :151--158; REVIEW:63--70,
:79--83; PROOF-COMPLETE].}  \textsc{Assume} Definition 90.H.  Let \(\alpha_L=\beta_{M_L,s_L}\circ\cdots\circ\beta_{M_1,s_1}\)
with \(M_i=L+i-1\), \(s_i=2i-1\); equivalently, with \(B_L:W_L\otimes W_L\to W_{2L}\) interleaving
\(\mathrm{old}_1..\mathrm{old}_L,\mathrm{new}_1..\mathrm{new}_L\) into \(\mathrm{old}_1,\mathrm{new}_1,\dots,\mathrm{old}_L,\mathrm{new}_L\)
by a crossing \(t_O\) of \(\mathrm{old}_j\) with \(\mathrm{new}_i\) exactly when \(i<j\),
\[
 \alpha_L(a)=B_L(a\otimes1_{W_L})B_L^\dagger . \tag{90.2}
\]
\textsc{Prove}: (i) \(\alpha_L\) is a unital injective \(*\)-homomorphism with \(\alpha_L(a)V_L=V_La\);
(ii) \(\alpha_L(\Alg_{[j,k]})\subset\Alg_{[2j-1,2k-1]}\), uniformly in \(L\), \(n\) and \(\Cc\), so \(r=0\) in the brief's
convention, and \(\alpha_L(f^{(j)})=f^{(2j-1)}\) for \(f\in\operatorname{End}(O)\); (iii) \([N_{2L},\alpha_L(a)]=\alpha_L([N_L,a])\);
(iv) for a unitary (braided, at \(\mathsf L1\)) strong monoidal equivalence with unitary \(u:F(O)\to O'\) carrying
\((F(\eta),F(t))\) to \((\eta',t')\), \(F_{2L}\alpha_L=\alpha'_LF_L\); (v) for \(L_n=2^nL_0\) the composites
\(\alpha^k_n=\alpha_{n-1}\cdots\alpha_k\) form a coherent directed system of unital injective \(*\)-maps (WT-1 A1).
\par\(\langle1\rangle1\). Before step \(i\), \(\mathrm{old}_i\) sits at \(2i-1\) and the suffix is \(\mathrm{old}_{i+1..L}\); so
step \(i\) inserts right after \(\mathrm{old}_i\), and expanding suffix crossings by the tensor identity gives (90.2).
\par\(\langle1\rangle2\). (i): Lemma 90.0 for each \(\beta\); the one-gap isometries compose to \(V_L\) (CA-71:102--113, (t)).
\par\(\langle1\rangle3\). (ii): by (90.1), insertions \(i<j\) shift both endpoints, \(j\le i<k\) only the right one, \(i\ge k\)
neither; hence the length-\((2(k-j+1)-1)\) interval
\[
 [\,j+(j-1),\ k+(j-1)+(k-j)\,]=[2j-1,2k-1]. \tag{90.3}
\]
\par\(\langle1\rangle4\). (iii): naturality gives \(t_O(p_e\otimes1)=(1\otimes p_e)t_O\); the \(\eta\) identity and unitarity give
\(t_Y(1_Y\otimes p_0)t_Y^\dagger=p_0\otimes1_Y\) (the one place pointedness enters).  Summing strand by strand,
\[
 N_{2L}B_L=B_L(N_L\otimes1+1\otimes N_L),\qquad [N_{2L},\alpha_L(a)]=\alpha_L([N_L,a]). \tag{90.4}
\]
\par\(\langle1\rangle5\). (iv): \(F\) and the tensorators (Penneys:325--348) carry the finite composite \(B_L\) to \(B'_L\).
\emph{Law 1 gap:} braided and unitary monoidal functors are not defined at a cited local line (REVIEW:30).  Compatibility
with composition of equivalences is asserted (NOTE:113), not proved, and not used (REVIEW:68).
\par\(\langle1\rangle6\). (v): cocycle by associativity of composition; WT-1 A1 and the norm inductive limit
(2010.11121:164--176).  Path independence across placement forests is not asserted (A1 does not need it).  \(\square\)

\emph{Consequences [derived].}  For \(h\in\Alg_{[j,j+1]}\), \(\alpha_L(h)\in\Alg_{[2j-1,2j+1]}\) and
\(\alpha_L(h)P_L=\theta_L(h)\): a \emph{homomorphic} local representative replacing the local representative of
Correction 76.B (INTEGRATION:80--82).  The lift transports constants with unit \(1\), the corner with unit \(P_L\); these
are different observables (NOTE:263).  \(\alpha_L\) is not reflection-equivariant (placement \(2j-1\), one crossing type;
the mirror is the opposite-crossing lift with even placement; INTEGRATION:69--71).

\textbf{Negative 90.N [derived, negative; NOTE:119; REVIEW:70].}  By (90.3), \(\alpha_L(N_L)=\sum_jp_1^{(2j-1)}\), so
\(N_{2L}-\alpha_L(N_L)=\sum_jp_1^{(2j)}\), of norm exactly \(L\) when \(X\ne0\).  Reading (iii) as a uniformly bounded
\emph{absolute-number} difference is false; what holds is the (am) endomorphism-degree statement (90.4).

\subsection{Wilson Consequences}

\textbf{Corollary 90.W [derived; E1.3, NOTE:123--134; REVIEW:72--73].}  \textsc{Assume} Theorem 90.1, states \(\omega_n\) on
\(\Alg_{L_n}\), and the (an) quantities with \(W_n=V_n\).  \textsc{Prove} \(d_{k,n}\le p_ne_n+2(1-p_n)\) for \(k\le n\) (\(p_ne_n:=0\) at \(p_n=0\)).
\(\langle1\rangle1\). Theorem 90.1(i) is the intertwining hypothesis of Proposition 76.1(iii)--(iv) (CA-76:90--108).
\(\langle1\rangle2\). At \(p_n=0\), \(\tau_n,e_n\) are undefined and the bound \(2\) is the trivial one; at \(p_n=1\) the
bound is \(e_n\), not \(0\).  \(\square\)  Summable bounds would give WT-1 A3 (\path{wilson_triangle.md}:18--107; computable
tails for WT-2, :109--129); none is supplied here.

\textbf{Proposition 90.V [derived; E1.4, NOTE:136--147; REVIEW:75--78].}  (a) \(e_M=\eta^{\otimes M}\) and
\(\omega^{\rm vac}_M(a)=e_M^\dagger ae_M\) satisfy \(e_{2M}=V_Me_M\), so \(\omega^{\rm vac}_{2M}\alpha_M=\omega^{\rm vac}_M\): \(p_n=1\),
\(e_n=0\), zero drift, and a WT-1 limit state exists.  This \(\mathsf R1\) instance is the \emph{trivial vacuum product
state}; no criticality is claimed.  (b) \emph{Negative.}  If \(X\ne0\), take a simple \(x\), an isometry \(u:x\to X\), and
\(f_M=(\iota u)\otimes\eta^{\otimes(M-1)}\); \(\omega^{\rm occ}_M(a)=f_M^\dagger af_M\) is also compatible.  Alternating
\(\omega_n\) between the two families gives \(p_n=1\) at every step, yet on \(2p_0^{(1)}-1\) adjacent pullbacks alternate
\(\pm1\): drift exactly \(2\), \(e_n=2\), no limit.  So \(p_n=1\) alone does not give A3.

\subsection{Non-Braided Input: Finite Equations and the One-Gap Obstruction}

Theorem 90.1 already runs at \(\mathsf L0\) with Definition 90.H; neither a braiding of \(\Cc\) nor a lift of \(O\) to
\(Z(\Cc)\) is needed (Remark 90.R).  The question is when the datum exists, and whether anything else can work.

\textbf{Proposition 90.2 [derived; (a) E1.6, NOTE:162--189, REVIEW:84--86; (b) reviewer Proposition R-1,
REVIEW:148--168].}  (a) Let \(s_i\) run over the simples of \(\mathcal D\) (unit included) and
\(f_{ij;k,\mu}:s_k\to s_i\otimes s_j\) be complete orthonormal isometric embeddings, all multiplicities kept.  Pointed
unitary relative half-braidings are exactly the solutions of the finite system (associators on both sides)
\[
 z_\one=1_O,\quad z_i(\eta\otimes1_i)=1_i\otimes\eta,\quad
 (f_{ij;k,\mu}\otimes1_O)z_k=(1_i\otimes z_j)(z_i\otimes1_j)(1_O\otimes f_{ij;k,\mu}),\quad z_i\text{ unitary}. \tag{90.5}
\]
A necessary integer condition is
\[
 \Delta_{x,c}:=\dim\Hom(c,O\otimes x)-\dim\Hom(c,x\otimes O)=0\ \ \forall\,x,c\ \text{simple},
 \quad\text{i.e. } [O][x]=[x][O]\ \text{in } K(\mathcal D). \tag{90.6}
\]
\(\langle1\rangle1\). Necessity: naturality, tensor identity, unitarity, centrality of \(\eta\); a unitary is an isomorphism in
each charge block, so \(\Delta\ne0\) excludes (90.5) (cf.\ ENO:1221).
\(\langle1\rangle2\). Sufficiency: \(z_Y=\sum_{i,\mu}(v_{i,\mu}\otimes1_O)z_i(1_O\otimes v_{i,\mu}^\dagger)\) is unitary and
basis-independent, natural by Schur, tensorial by summing (90.5) against \(f^\dagger\).  \(\square\)  (90.6) is only
necessary; the associator-dependent equations (90.5) remain.
(b) \textsc{Assume} \(M\ge2\), \(1\le s\le M-1\), \(J=W_s\), \(R=W_{M-s}\), and \emph{any} unital \(*\)-homomorphism
\(\beta:\Alg_M\to\Alg_{M+1}\) with \(\beta(a)v_{M,s}=v_{M,s}a\), \(\Alg_{[j,k]}\to\Alg_{[j,k]}\) for \(k\le s\) and
\(\Alg_{[j,k]}\to\Alg_{[j+1,k+1]}\) for \(j>s\) (straddling intervals unconstrained).  \textsc{Prove} \(\Delta_{x,c}=0\) for every
simple \(c\) and every simple \(x\) that is a summand of both \(J\) and \(R\).
\(\langle1\rangle1\). \(\beta(a\otimes1_R)=a\otimes1_O\otimes1_R\) and \(\beta(1_J\otimes b)=1_J\otimes1_O\otimes b\), by the
pattern, intertwining and \(v^\dagger v=1\).
\(\langle1\rangle2\). On \(S=\operatorname{End}(J)\otimes\operatorname{End}(R)\), block \((i,j)\) sits in block \(c\) of
\(\operatorname{End}(J\otimes R)\) with multiplicity \(N_{ij}^c\) and in block \(c'\) of \(\operatorname{End}(J\otimes O\otimes
R)\) with multiplicity \(T_{ij}^{c'}=\dim\Hom(c',s_i\otimes O\otimes s_j)\).
\(\langle1\rangle3\). With \(\Lambda\) the inclusion matrix of \(\beta\), \(\sum_cN_{ij}^c\Lambda_{cc'}=T_{ij}^{c'}\) (uniqueness
of decomposition of finite-dimensional representations).
\(\langle1\rangle4\). \(i=\one\) (as \(\one\subset O\subset J\)) gives \(\Lambda_{xc'}=\dim\Hom(c',O\otimes x)\); \(j=\one\) gives
\(\Lambda_{xc'}=\dim\Hom(c',x\otimes O)\).  \(\square\)
So (90.6) is necessary for \emph{every} one-gap lift with the support pattern (90.1), not only for the half-braiding
construction, once \(x\) occurs in \(W_s\) and \(W_{M-s}\).  \textbf{Not decided:} dyadic lifts with slack \(r\) or another
internal placement, where the left/right images are only lifts with unknown inclusion matrices (REVIEW:165--168).

\textbf{Boundary [open; NOTE:191, :269].}  (i)--(iii) constrain algebra maps, not crossings; \(z\) is not reconstructed
from an arbitrary \(\alpha\).  Existence or a complete obstruction for arbitrary local dyadic unital lifts at \(\mathsf L0\) is open.

\subsection{Circuit Depth and Dressing}

\textbf{Proposition 90.3 [derived; E1.7, NOTE:195--214; REVIEW:91--102].}  \textsc{Assume} Theorem 90.1, \(X\ne0\), the fine
\(2L\)-site chain and the contiguous input \(a\mapsto a\otimes1_{W_L}\).  \textsc{Prove}: \(B_L\) is a nearest-neighbour
circuit of \(t_O\) crossings of depth \(L-1\) with \(L(L-1)/2\) gates, and any range-\(w\), depth-\(D\) circuit \(U\) with
\(U(a\otimes1)U^\dagger=\alpha_L(a)\) obeys \(D(w-1)\ge L-1\) \hfill\textup{(90.7)}
\(\langle1\rangle1\). Layer \(t=1,\dots,L-1\): crossings on \((p,p+1)\), \(p=L-t+2i-1\), \(i=1,\dots,t\), exchanging
\(\mathrm{old}_{L+i-t}\) with \(\mathrm{new}_i\); disjoint and adjacent by induction on \(t\).
\(\langle1\rangle2\). Each pair \(i<j\) crosses once (at \(t=L+i-j\)) and the final word is interleaved.  A shuffle of two
increasing sequences is \(321\)-avoiding, so all its reduced words are commutation-equivalent: no braid
(Yang--Baxter) move, which a relative half-braiding need not supply among old strands, is used, and tensor interchange
identifies the schedule with \(B_L\) (the combinatorial lemma made explicit at REVIEW:92--97).
\(\langle1\rangle3\). One layer enlarges an interval support by at most \(w-1\) at each end (interchange).
\(\langle1\rangle4\). \(p_0^{(L)}\otimes1\) has support \(\{L\}\) and image \(p_0^{(2L-1)}\) (Theorem 90.1(ii)).  A \(p_0\) outside an
interval is not in that interval algebra: compressing the other sites to \(\eta\), interval operators give equal scalars on
\(\eta\) and on an isometric simple summand of \(X\) at the omitted site, \(p_0\) gives \(1\) and \(0\).  \(\square\)
So no fixed depth and range implement the family in this geometry, while \(\alpha_L\) itself has zero support growth.
Distributed-ancilla geometries change the input convention and are not compared (NOTE:216).

\textbf{Corollary 90.D [derived implication; NOTE:218--224; REVIEW:104].}  For a \emph{supplied} output circuit \(U\) of
depth \(D\) and range \(w\), \(W=UV_L\) and \(\alpha^U=\mathrm{Ad}_U\circ\alpha_L\) satisfy \(\alpha^U(a)W=Wa\), and an \(m\)-site
input has output support length \(\le2m-1+2D(w-1)\): \(r=2D(w-1)\).  Degree preservation needs \([U,N_{2L}]=0\) in addition.
\(U\) is not constructed; a quasi-local \(U\) needs its own tail estimate (D3's dressing lemma stays open).

\subsection{Integration Status and the Next Lemma}

CA-75 Link (b1) is closed at \((\mathsf{any};\ \mathsf L0+\text{H};\ \mathsf R0^u)\); without the datum at \(\mathsf L0\) it is
open.  Proposal 76.5 is superseded by Theorem 90.1, including the owed non-pointed check with sourced \(F,R\)
(Fibonacci row).  D3-3's local-witness route (Question 76.12 (L)) becomes applicable; its witness hypotheses remain.
Corrections 76.A, 76.C and Negative 76.11(d) are untouched.  Theorem 90.1 is open-chain only: periodic seam coherence of
\(B_L\) (needed by \(\mathsf S_{\rm full}\), LC-4) is open (NOTE:273).

\textbf{Target 90.T [proposal; dyadic transport-extraction; INTEGRATION:85--91].}  Let \(\alpha_L:\Alg_L\to\Alg_{2L}\) be
unital \(*\)-maps with \(\alpha_L(a)V_L=V_La\) and \(\alpha_L(\Alg_{[j,k]})\subset\Alg_{[2j-1-r,\,2k-1+r]}\) for a fixed \(r\) and
all \(L\).  Show that the Bratteli argument of 90.2(b) survives the slack (block \(r\) sites at the cut, track the inclusion
matrices of the left and right partial lifts), forcing \([O][x]=[x][O]\) in \(K(\mathcal D)\) for every simple \(x\) of
\(\mathcal D\); or exhibit a local dyadic lift for the \(\mathrm{Vec}_{S_3}\) row.  Either outcome settles \(\mathsf R0^u\) at
\(\mathsf L0\) up to the associator equations (90.5).

\subsection{What This Does Not Claim}

No WT-1 A3 for any non-trivial state family; \(\mathsf R1\) is reached only for the vacuum product state.  No critical
state, gap, KS convergence, or rung \(\ge\mathsf R2\); no periodic or parity-covariant lift; no decision for arbitrary local
dyadic lifts at \(\mathsf L0\), and (90.6) is not claimed sufficient.  Naturality (iv) concerns the tower construction, not
\(\mathsf R5\).  The depth bound concerns contiguous ancillas only.  The finite checks are Float64 calibrations; the
``not braided'' status of \(\mathrm{Vec}_{\mathbb Z_3}^\omega\) is a finite-grid search, not a proof (REVIEW:189--190).

\subsection{Examples (Instances Only)}

\begin{center}\footnotesize
\begin{tabular}{p{3.3cm}p{10.7cm}}
\toprule
Instance & Calibration datum for the general statements \\
\midrule
any \(\Cc\), \(X=0\); \(\mathrm{Vec}\), \(O=\mathbb C^d\) & \(\alpha_L=\mathrm{id}\) (excluded from (90.7)); tensor shuffle with \(t=\) flip (REVIEW:144). \\
Fibonacci, \(O=\one\oplus\tau\) & Sourced \(F,R\) (\path{RSWfinal3.tex}:2379--2407), \(\le4\) sites: all matrix-unit tests pass; phase-free crossing fails block naturality, error \(0.786\) (\path{output.txt}:4--6, \path{red_phase_free.txt}:12). \\
\(\mathrm{Vec}_{\mathbb Z_3}\), \(O=\one\oplus g\) & Non-self-dual \(X\): passes \(L=2,3\); the \(g\)--\(g\) phase is tested only at \(L=4\) (\(\mathbb Z_n\) needs \(n\) old sites) (\path{review_pointed_lift_output.txt}:2--3). \\
semion, \(O=\one\oplus2g\); \(\one\oplus\one'\oplus g\) & Multiplicity, nontrivial associator, unit summand in \(X\): pass; phase-free crossing fails support (:4--6). \\
\(\mathrm{Vec}_{\mathbb Z_3}^\omega\), \(O=\one\oplus g\) & Non-braided (grid search): three pointed half-braidings; Theorem 90.1 checks pass at \(L=4\); also \(O=\one\oplus g\oplus g^2\) (:7--11). \\
\(\mathrm{Vec}_G\), \(O=\one\oplus g\), \(g\) noncentral & \(\langle g\rangle\) cyclic with symmetric braiding, so 90.1 applies; \([O][h]-[h][O]=[gh]-[hg]\ne0\): no lift to \(Z(\mathrm{Vec}_G)\) (NOTE:283). \\
\(\mathrm{Vec}_{\mathbb Z/2}\), non-central point & Half-braiding on \(O\) with \(\eta\) not central: \(\alpha(h)V=-Vh\) (NOTE:284; \path{review_pointed_lift_output.txt}:12). \\
\(\mathrm{Vec}_{S_3}\), \(O=\one\oplus g\oplus h\), \(gh\ne hg\) & \(\Delta\ne0\) on \(\langle O\rangle=S_3\) (four clashing entries at \(M=2,s=1\)): no one-gap lift (90.2(b)); dyadic lifts undecided (Target 90.T) (\path{review_one_gap_obstruction_output.txt}:1--4). \\
\bottomrule
\end{tabular}
\end{center}
